% 87-02 ③(87쪽) — 보기 그래프: $x=-1$ 에서 극소, $x=3$ 에서 기울기 0 인 아래로 볼록한 사차함수 개형. 그림 자체가 보기라 TikZ 로 다시 그림.
% 라벨은 원문 그대로 -1, O, 3, x, y, y=f(x) 와 점선 두 줄만.
\begin{tikzpicture}[baseline=(current bounding box.center), x=0.36cm, y=0.36cm, line width=0.5pt]
  \path (-2.9,-2.4) rectangle (5.15,3.5);
  \draw[->, >=stealth, line width=0.7pt] (-2.7,0) -- (4.8,0) node[below=1pt, font=\footnotesize] {$x$};
  \draw[->, >=stealth, line width=0.7pt] (0,-2.0) -- (0,3.3) node[left=1pt, font=\footnotesize] {$y$};
  \draw[densely dotted, line width=0.4pt] (-1,0) -- (-1,-1.15);
  \draw[densely dotted, line width=0.4pt] (3,0) -- (3,1.05);
  \draw[line width=0.6pt, smooth] plot coordinates {(-2.4,2.9) (-2.25,2.2) (-2.1,1.5) (-1.95,0.9) (-1.8,0.4) (-1.6,-0.15) (-1.4,-0.65) (-1.2,-1.0) (-1,-1.15) (-0.8,-1.12) (-0.6,-1.05) (-0.3,-0.9) (0,-0.72) (0.35,-0.5) (0.7,-0.25) (0.95,-0.05) (1.3,0.2) (1.7,0.45) (2.1,0.7) (2.5,0.9) (2.8,1.0) (3,1.05) (3.3,1.15) (3.6,1.45) (3.85,1.95) (4.1,2.7)};
  \node[above=1pt, font=\footnotesize] at (-1.05,0) {$-1$};
  \node[below left=-2pt and -3pt, font=\footnotesize] at (0,0) {$\pt{O}$};
  \node[below=1pt, font=\footnotesize] at (3,0) {$3$};
  \node[anchor=west, font=\footnotesize] at (0.25,3.05) {$y=f(x)$};
\end{tikzpicture}
